On considère les fonctions $f_{k}$ définies sur $\R$ par 
$f_{k}(x) = x + k \e^{-x},$
 où $k$ est un réel strictement positif.

	\begin{enumerate}
		\item On s'intéresse dans cette question au cas $k = 0,5$, donc à la fonction $f_{0,5}$ définie sur $\R$ par : $f_{0,5}(x) = x + 0,5 \e^{-x}.$

\smallskip
	\begin{enumerate}
		\item 
$f'_{0,5}(x)=1+0,5\times(-1)\e^{-x}=1-0,5\e^{-x}$		
		
		\item 
$\begin{gathered}f'_{0,5}(x) >0
\iff
1-0,5\e^{-x} >0
\iff
1>0,5\e^{-x}
\iff
\dfrac{1}{0,5}>\e^{-x}\\
\phantom{f'_{0,5}(x) >0}
\iff
2>\e^{-x}
\iff
\ln(2)>-x
\iff
-\ln(2)<x
\iff
x>\ln\left (\dfrac{1}{2}\right )\\
\phantom{f'_{0,5}(x) >0}
\iff
x>\ln(0,5)\end{gathered}$

De même: $f'_{0,5}(x)=0 \iff x=\ln(0,5)$, et $f'_{0,5}(x)<0 \iff x<\ln(0,5)$.

On en déduit que la fonction $f_{0,5}$ admet un minimum en $\ln(0,5)$.
	\end{enumerate}
\end{enumerate}

Soit $k$ un réel strictement positif. On donne le tableau de variations de la fonction $f_{k}$.
\medskip

\begin{center}
{\renewcommand{\arraystretch}{1.3}
\psset{nodesep=3pt,arrowsize=2pt 3}  
\def\esp{\hspace*{1.5cm}}
\def\hauteur{0pt}
$\begin{array}{|c| *4{c} c|}
\hline
 \text{Valeurs de } x & -\infty & \esp & \ln(k) & \esp & +\infty \\
\hline
  & \Rnode{max1}{+\infty}  &  &  &  & \Rnode{max2}{+\infty}   \\
\text{Variations de } f_k & &  & & &  \rule{0pt}{\hauteur} \\
 &  & &   \Rnode{min}{f_k(\ln k)} & & \rule{0pt}{\hauteur}
\ncline{->}{max1}{min} \ncline{->}{min}{max2}\\
\hline
\end{array}$
\renewcommand{\arraystretch}{1}}
\end{center}

	\begin{enumerate}
		\setcounter{enumi}{1}
		\item 
$f_k(\ln k) = \ln(k)+ k\e^{-\ln(k)} = \ln(k) + \dfrac{k}{\e^{\ln(k)}} = \ln(k) + \dfrac{k}{k} = \ln(k) + 1$


	\end{enumerate}

On note $\mathcal{C}_{k}$ la courbe représentative de la fonction $f_{k}$ dans un plan muni d'un repère orthonormé.

On note $\mathrm{A}_{k}$ le point de la courbe $\mathcal{C}_{k}$ d'abscisse $\ln k$.

On a représenté ci-dessous quelques courbes $\mathcal{C}_{k}$ pour différentes valeurs de $k$.

\begin{center}
\psset{unit=1cm,algebraic}
\def\xmin{-4.2} \def\xmax{9} \def\ymin{-1} \def\ymax{9}
\begin{pspicture*}(\xmin,\ymin)(\xmax,\ymax)

\psaxes[arrowsize=3pt 3, ticksize=-2pt 2pt]{->}(0,0)(\xmin,0)(\xmax,\ymax) 

\multido{\n=0.5+0.5}{10}{
\psplot[plotpoints=1000,linecolor=red]{\xmin}{\xmax}{x+\n*EXP(-x)}
}
\listplot[plotstyle=dots]{-0.693	0.307 0.000	1.000 0.405	1.405 0.693	1.693 0.916	1.916 1.099	2.099 1.253	2.253 1.386	2.386 1.504	2.504 1.609	2.609}
{\footnotesize
\uput[u](-0.693,0.307){A$_{0,5}$} \uput[r](0,1){A$_{1}$} 
\uput[r](0.405,1.405){{A$_{1,5}$}}  \uput[u](1.609,2.609){A$_{5}$}}
\psplot[linecolor=blue,linestyle=dashed]{-1}{5}{x+1}
\uput[u]{45}(3,4){\blue $y=x+1$}
\end{pspicture*}
\end{center}

\begin{enumerate}
		\setcounter{enumi}{2}
		\item 
Le point A$_k$ a pour coordonnées $(\ln(k)\;;\; f_k(\ln(k)))$ soit $(\ln(k)\;;\; \ln(k)+1)$. Donc tous les points A$_k$ sont situés sur la droite d'équation $y=x+1$.

L'affirmation : \og{}Pour tout réel $k$ strictement positif, les points A$_{0,5}$, A$_{1}$ et A$_{k}$ sont alignés. \fg{} est donc vraie.

\end{enumerate}
